\documentclass[11pt,a4paper]{article}
\usepackage[margin=1in]{geometry}
\usepackage[T1]{fontenc}
\usepackage[utf8]{inputenc}
\usepackage{lmodern}
\usepackage{hyperref}
\usepackage{enumitem}
\usepackage{booktabs}

\title{Scriptura API: Technical Report}
\author{Hansel D'Cruz}
\date{11 March 2026}

\begin{document}

\maketitle

\section{Project Overview}
Scriptura API is a REST API for the Douay-Rheims Bible corpus. It provides hierarchical text access (books, chapters, verses), user-curated collections, theme coverage, and lexical similarity recommendations.

The design goal was to meet CRUD requirements and add useful analytics without overengineering. The API returns JSON over HTTP and is documented with OpenAPI/Swagger.

\section{Technology Stack and Rationale}
\subsection{Programming Language: Python}
I selected Python because it allows for quick prototyping, is easy to read, and has a lower learning curve for this assessment. The priority here was API quality, not spending time learning extra tooling.

In this project, Python fits naturally because I perform statistical and text-analysis tasks for theme coverage and lexical similarity. It let me implement both standard REST endpoints and analytics endpoints in one consistent codebase. Python's concise syntax made it practical to quickly build and refine collection CRUD, theme coverage analytics, similarity graph generation, and verse recommendation flows. I also have industry experience with Python in web and data-driven work, so I could apply good practices early, including project structure, validation, and tests for critical paths. The language is clear to review, especially in analytics services and serializer logic, and the ecosystem used in this project (Django, DRF, django-filter, drf-spectacular, django-allauth) reduced low-level boilerplate and potential implementation issues.

\subsection{Framework: Django + Django REST Framework}
I chose Django 4.2 with Django REST Framework (DRF) for structure, speed, and reliability.

Django follows the Model-Template-View (MTV) pattern, conceptually close to MVC. In this API focused project, this nicely creates a practical separation of concerns:
\begin{itemize}[leftmargin=*]
	\item \textbf{Models}: Represent scripture hierarchy (Book, Chapter, Verse) and user curation entities (Collection, Theme relationships).
	\item \textbf{Views/ViewSets}: Handle HTTP request-response lifecycle, including standard resource endpoints and custom analytics endpoints.
	\item \textbf{Serializers}: Define API contracts, validation, and transformation between model instances and JSON payloads.
\end{itemize}

Additional reasons for Django/DRF selection:
\begin{itemize}[leftmargin=*]
	\item \textbf{Integrated features}: ORM, migrations, authentication, admin, and secure defaults reduced boilerplate.
	\item \textbf{API productivity}: DRF ViewSets, routers, pagination, and filtering kept endpoint design consistent.
	\item \textbf{Documentation}: \texttt{drf-spectacular} generated OpenAPI docs that are easy to test in Swagger.
	\item \textbf{Maintainability}: Django app boundaries align with this codebase split (core, analytics, themes, ingestion).
\end{itemize}

\subsection{Database: SQLite and Why}
The current implementation uses SQLite for local development and demonstration. This fits the project because the data is strongly relational: \texttt{Book $\rightarrow$ Chapter $\rightarrow$ Verse}, plus linked entities such as \texttt{Section}, \texttt{Footnote}, \texttt{Theme}, and \texttt{Collection}. These relationships are the bulk of my API and are queried frequently through filtered lookups.

An SQL database is almost definitely a better fit than a NoSQL document store for this model because it gives explicit foreign keys, referential integrity, many-to-many support, and transactional updates. In this project, that matters for collection updates (verses/themes links), analytics precomputation tables, and consistent retrieval of nested scripture data. A NoSQL design could duplicate data across documents and make relationship updates and consistency checks harder.

SQLite was selected for now because it has minimal setup (as it comes out of the box with Django), and a predictable local environment. The trade-off is concurrency at production scale. For future deployment, PostgreSQL would be preferred for stronger indexing, better concurrent write handling, and broader scaling options. Given the scale of this coursework however, this wasn't a meaningful consideration for me.

\subsection{Supporting Libraries}
Key supporting choices:
\begin{itemize}[leftmargin=*]
	\item \texttt{django-filter}: Query parameter filtering,
	\item \texttt{drf-spectacular}: OpenAPI schema and Swagger UI,
	\item \texttt{django-allauth}: Authentication integration (including Google OAuth),
	\item \texttt{python-dotenv}: Environment-variable based configuration,
	\item \texttt{django-cors-headers}: CORS management for browser-based clients.
\end{itemize}

\section{Architecture and Design Choices}
\subsection{Modular App Structure}
The project is split into focused Django apps:
\begin{itemize}[leftmargin=*]
	\item \textbf{core}: canonical text models and main API endpoints,
	\item \textbf{analytics}: text metrics, similarity graph, recommendation logic,
	\item \textbf{themes}: thematic classification and keyword management,
	\item \textbf{ingestion}: data-loading pipeline from USFM and related files.
\end{itemize}

\subsection{Layered Design Decisions}
I deliberately used a semi-layered architecture:
\begin{itemize}[leftmargin=*]
	\item \textbf{Repository layer} (e.g., \texttt{core/repositories.py}) isolates ORM query access,
	\item \textbf{Service layer} (analytics services) holds reusable NLP/statistical logic,
	\item \textbf{Serializer split} between read and write serializers for \texttt{Collection} to control response shape and validation,
	\item \textbf{ViewSets + APIViews}: ViewSets for standard CRUD/read routes, APIViews for custom analytics workflows.
\end{itemize}

\subsection{Data Model and Relationships}
The data model supports both scripture structure and user interaction:
\begin{itemize}[leftmargin=*]
	\item \texttt{Book $\rightarrow$ Chapter $\rightarrow$ Verse} hierarchical traversal,
	\item Optional \texttt{Section} and \texttt{Footnote} for richer text context,
	\item \texttt{Collection} as user-curated many-to-many mapping to verses and themes,
	\item \texttt{BookSummary}, \texttt{ThemeCoverageCache}, and \texttt{SimilarityCache} for precomputed/cached analytics payloads.
\end{itemize}

\subsection{API Design and Error Conventions}
The API uses REST-style endpoint naming and standard status codes:
\begin{itemize}[leftmargin=*]
	\item \texttt{200 OK} for successful reads/updates,
	\item \texttt{201 Created} for successful resource creation,
	\item \texttt{400 Bad Request} for invalid query/body input,
	\item \texttt{404 Not Found} for missing resources.
\end{itemize}
All endpoints return JSON for consistent client behaviour.

\section{Testing Strategy and Quality Approach}
Testing was implemented with Django's test framework and the DRF test client.

\subsection{What Was Tested}
\begin{itemize}[leftmargin=*]
	\item Unit checks for analytics metrics (word count, entropy, TTR, hapax),
	\item API smoke tests for core routes and filtering/search behaviour,
	\item Cache behaviour tests for theme coverage (creation and refresh on keyword changes),
	\item Endpoint-level checks for status-code correctness.
\end{itemize}

\subsection{Error Handling and Robustness}
Implemented edge handling includes:
\begin{itemize}[leftmargin=*]
	\item Validation of similarity metric values,
	\item Threshold parsing/clamping for numeric query parameters,
	\item Defensive checks for required parameters like \texttt{verse\_id} and \texttt{collection\_id},
	\item Not-found behaviour for missing verses/collections/themes.
\end{itemize}

\section{Challenges, Reflection, and Lessons Learned}
\subsection{Main Challenges}
The biggest challenge was data sourcing and normalisation, not endpoint coding. I initially planned to use the KJV text, but I moved away from it due to licensing constraints (including UK Crown rights considerations). I then switched to an open source Douay-Rheims American (DRA) source.

At first, I used public JSON datasets from GitHub. They were easy to ingest, but the reading experience felt unnatural compared to reading an actual Bible. I realised that what was missing were not just verses, but the structure as well. Pericope headers, footnotes, and paragraph-level segmentation are present within the actual text. Most JSON sources I found did not carry enough metadata for these features.

To address this, I used AI to discover potential solutions and ended up discovering USFM, which is a more standardised scripture format (\url{https://ubsicap.github.io/usfm/}), and used the public DRA text distribution from eBible (\url{https://ebible.org/find/show.php?id=engDRA}). This resolved most paragraph and structural rendering issues. For pericope groupings, I used a public-domain dataset from Hugging Face (\url{https://huggingface.co/datasets/JWBickel/KJV_Pericopes}) and mapped headings to verses where possible.

There were still gaps: modern KJV-oriented pericope resources usually exclude deuterocanonical books (since Protestants do not consider them canon for historical reasons), so those books have limited header coverage. Given the time constraints, I did not build a separate pericope pipeline for the deuterocanonical set - but this is something to be added in the future.

For footnotes, I worked from the Project Gutenberg DRA text (\url{https://www.gutenberg.org/cache/epub/8300/pg8300.txt}), combined AI-assisted parsing with manual clean-up, and integrated the result into the ingestion pipeline. A practical issue here was name mapping, because Old and New Testament book names vary between sources and there is no single universal naming convention.

Beyond data integration, there were standard engineering challenges. Nested scripture responses can be expensive, pairwise similarity over books is computationally heavy, and collection updates require careful many-to-many validation.

\subsection{Lessons Learned}
\begin{itemize}[leftmargin=*]
	\item Data engineering decisions can dominate API projects; format and metadata quality matter as much as schema design.
	\item Open data is rarely unified, so ingestion must handle naming mismatches, missing fields, and editorial differences across sources.
	\item Early architectural boundaries (repositories/services) reduce refactor cost when data pipelines evolve.
	\item Cache key strategy matters as much as cache storage itself for analytics endpoints.
\end{itemize}

\section{Limitations and Future Improvements}
\subsection{Current Limitations}
\begin{itemize}[leftmargin=*]
	\item SQLite limits production-grade concurrency.
	\item Collection access policy is intentionally permissive in places and can be tightened.
	\item Recommendation quality is lexical and bag-of-words based; deeper semantic/contextual methods are not yet used.
	\item Test coverage is present but can be expanded for edge cases and permissions.
\end{itemize}

\subsection{Future Development}
\begin{itemize}[leftmargin=*]
	\item Migrate to PostgreSQL and add tuned indexes,
	\item Strengthen access control and object-level permissions,
	\item Add async/background task processing for heavy analytics,
	\item Add benchmark suite and load testing,
	\item Evaluate embedding-based semantic retrieval alongside TF-IDF/cosine pipelines,
	\item Package CI checks (tests + lint + schema validation) in GitHub Actions.
\end{itemize}

\section{Generative AI Declaration and Critical Analysis}
\subsection*{Declaration}
Generative AI was used extensively across the project, with manual changes made where required by my own design preferences and implementation constraints. In practice, this included boilerplate generation, UI support (including HTML/CSS components), ingestion script drafting, model/serializer scaffolding, and initial unit-test generation. Final design decisions, integration, debugging, and acceptance of all outputs were performed by me.

\subsection*{How GenAI Was Used Responsibly}
\begin{itemize}[leftmargin=*]
	\item \textbf{Boilerplate and UI}: GenAI was used to speed up repetitive implementation tasks. An example prompt was: \textit{``Generate a search bar on the homepage with state management and dynamic autosuggest using my database as the source.''}
	\item \textbf{Data ingestion support}: I manually researched data sources and formats, then used GenAI to draft ingestion scripts after I provided source structure details and edge cases.
	\item \textbf{Schema-driven backend generation}: I provided the database schema and relationship structure (Books, Chapters, Verses, plus one-to-one and many-to-many links), then used GenAI to generate draft models, serializers, and tests.
	\item \textbf{Unit tests}: I specified what needed to be tested and then verified that generated tests were valid and meaningful.
	\item \textbf{Manual validation}: Outputs were treated as proposals, reviewed line by line, and edited where I needed it to be.
\end{itemize}

\subsection*{Critical Reflection on GenAI Usage}
\begin{itemize}[leftmargin=*]
	\item \textbf{Benefits}: It improved the speed of the project and reduced the overhead of learning syntax.
	\item \textbf{Risks}: Inaccurate assumptions, over-generalised patterns, and hidden logic errors if accepted uncritically.
	\item \textbf{Project-specific risk observed}: Some generated outputs assumed the Protestant canon (66 books compared to Catholic canon containing 73) rather than the Catholic canon used here, so I had to manually correct book coverage and related assumptions.
	\item \textbf{Mitigation}: Explicit verification with tests, runtime checks, and source-level review before committing any AI-assisted change.
\end{itemize}

\end{document}
